\chapter{Legal swaps and permutohedral complexes}\label{ch:swaps}

This chapter separates what is true for \emph{every} set of orderings from what is true for arithmetically (prefix-)defined sets and holds only in a verified range. The universal statements are the reconstruction of $K_A$ from the generator-coloured swap graph (\cref{thm:swap-reconstruction}) and the completeness of face recognition (\cref{thm:face-recognition}(a)); the bounded ones are the graph-only recovery of cells (the six-cycle classification for $k\le8$ and the soundness checks).

\section{Squares, hexagons and generator colours}
The permutohedron graph is the Cayley graph of $S_k$ for $s_1,\dots,s_{k-1}$, where $x\mapsto xs_i$ exchanges the entries in positions $i,i{+}1$; the edge $\{x,xs_i\}$ has \emph{colour} $i$. The Coxeter relations are $s_i^2=e$, $(s_is_j)^2=e$ for $|i-j|\ge2$ (a commuting pair), $(s_is_{i+1})^3=e$ (a braid relation). They give the two kinds of $2$-faces: a \emph{square} $x\langle s_i,s_j\rangle$ ($|i-j|\ge2$, four vertices, boundary a $4$-cycle of alternating colours $i,j$) and a \emph{hexagon} $x\langle s_i,s_{i+1}\rangle$ ($6$ vertices, a $6$-cycle). For $J\subseteq\{1,\dots,k-1\}$ let $W_J=\langle s_j:j\in J\rangle$, a product of symmetric groups on the maximal runs of consecutive positions.

\begin{lemma}[faces are parabolic cosets]\label{lem:parabolic}
The faces of $P_k$ are in bijection with pairs $(J,xW_J)$: the face $(B_1|\dots|B_m)$ corresponds to $J=$ the set of positions $i$ such that positions $i,i{+}1$ lie in the same block, and to the coset $xW_J$ where $x$ is any vertex of the face. The vertex set of the face is $xW_J$. The subgraph of the permutohedron graph spanned by the colours in $J$ has the cosets $xW_J$ as its connected components.
\end{lemma}
\begin{proof}
A vertex of $(B_1|\dots|B_m)$ lists the blocks in order with arbitrary internal orders, i.e. $x$ composed with permutations of positions within each block, which is $W_J$. A $J$-coloured edge exchanges two positions inside a block, so never leaves a coset; the Cayley graph of $W_J$ for its Coxeter generators is connected.
\end{proof}

\section{Reconstruction from the coloured swap graph}
\begin{theorem}[reconstruction of $K_A$]\label{thm:swap-reconstruction}
Let $A\subseteq S_k$ be \emph{arbitrary} and let $G_A$ be its legal-swap graph with the colouring. Then the faces of $K_A$ are the pairs $(J,Z)$ with $Z$ a connected component of the $J$-coloured subgraph of $G_A$ having exactly $|W_J|$ vertices; face incidence is vertex-set inclusion with $|J|$ dropping by $1$. Hence the $f$-vector, all boundary matrices over $\F$ and all $\beta_j$ are functions of $(G_A,\text{colours},k)$. A colour-preserving isomorphism $G_A\to G_B$ up to a Coxeter-diagram automorphism ($\sigma=\mathrm{id}$ or $j\mapsto k-j$) induces an isomorphism of chain complexes, hence equal Betti numbers.
\status{\Analytic (no arithmetic hypothesis)}{GIG-I v002, T13}
\end{theorem}
\begin{proof}
\emph{Faces.} By \cref{lem:parabolic}, $xW_J\subseteq A$ iff the connected component $Z$ of $x$ in the $J$-coloured subgraph of $G_A$ has $|W_J|$ elements: ($\Rightarrow$) the whole coset lies in $A$ and is $J$-connected, and a component can never exceed the coset, so $Z=xW_J$; ($\Leftarrow$) $Z\subseteq xW_J$ always, so $|Z|=|W_J|$ forces $Z=xW_J\subseteq A$. \emph{Incidence.} For faces $F,F'$ of a polytope, $V(F')\subseteq V(F)$ iff $F'$ is a face of $F$; dimension drops by one iff $|J'|=|J|-1$. A regular cell complex has all incidence coefficients $1$ in codimension one, so the $\F$-boundary matrices are determined. \emph{Betti numbers.} $\beta_j=f_j-\rho_j-\rho_{j+1}$. \emph{Isomorphism.} $\sigma$ maps runs of consecutive colours to runs of equal length, so $|W_{\sigma(J)}|=|W_J|$, and the isomorphism maps $J$-components onto $\sigma(J)$-components preserving size and inclusion.
\end{proof}
\begin{proposition}[isomorphism of graphs preserves accessibility data]\label{prop:iso-profiles}
If $G_A\cong G_B$ as (uncoloured) graphs then component-size multisets, multisets of ball-size vectors $(|B_0(w)|,\dots,|B_r(w)|)$ for every $r$, the integers $P_r=\prod_w|B_r(w)|$, the profile $I_r=\log_2P_r/N$, class entropy and separating radius coincide.
\status{\Analytic}{GIG-I v002, T14}
\end{proposition}
\begin{proof}
These are defined by the graph metric alone; an isomorphism maps balls bijectively onto balls.
\end{proof}

\section{What the uncoloured graph determines}\label{sec:graph-only}
Without colours, additional structure is needed. For \emph{every} $k$ we have:
\begin{lemma}[four-cycles are squares]\label{lem:four-cycles}
Every $4$-cycle of the permutohedron graph is the boundary of a square: its colours alternate $a,b,a,b$ with $|a-b|\ge2$.
\end{lemma}
\begin{proof}
A $4$-cycle gives a word $s_as_bs_cs_d=e$ with $a\ne b$, $b\ne c$, $c\ne d$, $d\ne a$ (no backtracking, also cyclically). Case $a=c$: $s_as_bs_a=s_d$ is an adjacent transposition; if $|a-b|=1$ then $s_as_bs_a$ is the transposition of two positions at distance $2$, not adjacent; so $|a-b|\ge2$, $s_as_bs_a=s_b$, $d=b$. Case $a\ne c$: $s_as_b=s_ds_c$. If $|a-b|\ge2$ the left side is a product of two disjoint transpositions whose factorisations $s_xs_y$ are $(x,y)\in\{(a,b),(b,a)\}$; $(d,c)=(a,b)$ gives $c=b$, and $(d,c)=(b,a)$ gives $c=a$, both excluded. If $|a-b|=1$ the left side is a $3$-cycle whose only factorisation of the form $s_xs_y$ is $(a,b)$ (the other order is the inverse), giving $c=b$, excluded.
\end{proof}
Six-cycles require a classification that is currently available only for a finite range.
\begin{lemma}[six-cycles of the permutohedron graph, $k\le8$]\label{lem:six-cycles}
For $4\le k\le8$ every $6$-cycle of the permutohedron graph is of exactly one of four types: a \emph{braid hexagon} (a $2$-face of type $(3)$; chordless), a \emph{Petrie cycle} (six vertices of a commuting cube omitting two antipodal vertices; chordless), or a cycle with exactly one chord (a ``domino'' of two squares, or a mixed type). For a braid hexagon no vertex outside it is adjacent to two of its vertices; for a Petrie cycle each of the two omitted antipodes is adjacent to three of its vertices.
\status{\CompAss: enumeration of all $6$-cycles through a root vertex. Re-derived by the manuscript verifier for $k=4,5,6$ (through the root: $k=4$: $2$ braid; $k=5$: $3$ braid, $6$ chorded; $k=6$: $4$ braid, $3$ Petrie, $27$ chorded; \texttt{ch5\_six\_cycles\_through\_root\_k4/5/6}); $k=7,8$ inherited and not re-run}{GIG-I v002, E5 enumeration; T15 step 2}
\end{lemma}
The adjacency claim for braid hexagons holds for all $k$:
\begin{lemma}\label{lem:braid-isolated}
No vertex $u$ of $P_k$ outside a braid hexagon $H=xW$, $W=\langle s_i,s_{i+1}\rangle$, is adjacent to two vertices of $H$.
\end{lemma}
\begin{proof}
Suppose $u=x's_c=y's_d$ with $x'\ne y'$ in $H$. Then $c\ne d$ and $x'^{-1}y'=s_cs_d\in W$; moreover $s_c\notin W$ because $u\notin H$, i.e. $c\notin\{i,i+1\}$. Elements of $W$ are supported on the positions $\{i,i{+}1,i{+}2\}$. If $|c-d|\ge2$ the support of $s_cs_d$ has four positions; if $|c-d|=1$ it is the three positions $\{\min(c,d),\dots,\min(c,d)+2\}$, which equals $\{i,i+1,i+2\}$ only if $\{c,d\}=\{i,i+1\}$, contradicting $c\notin\{i,i+1\}$.
\end{proof}

\begin{theorem}[two-skeleton from the uncoloured graph]\label{thm:skeleton}
Let $A\subseteq S_k$ have the cube-product property (CP), in particular any prefix-defined $A$ (\cref{cor:CP}), and let $4\le k\le8$. Then the $2$-faces of $K_A$ are exactly the $4$-cycles of the \emph{uncoloured} graph $G_A$ and the chordless $6$-cycles of $G_A$ having no vertex of $G_A$ outside the cycle adjacent to two of its vertices. Consequently $f_0,f_1,f_2$, $\rho_1,\rho_2$, $\beta_0$ and $\beta_1$ of $K_A$ are functions of the isomorphism class of the uncoloured $G_A$.
\status{\Analytic given \cref{lem:six-cycles} (computer-assisted, $k\le8$)}{GIG-I v002, T15}
\end{theorem}
\begin{proof}
\emph{Squares.} By \cref{lem:four-cycles} every $4$-cycle of $G_A$ bounds a square of $P_k$ whose four vertices are in $A$, hence a face of $K_A$; every square face gives a $4$-cycle. \emph{Hexagons.} By \cref{lem:six-cycles} a chordless $6$-cycle in $G_A$ is a braid hexagon or a Petrie cycle (chords of $G_A$ are chords in $P_k$ because all cycle vertices are in $A$). For a braid hexagon, by \cref{lem:braid-isolated} no vertex is adjacent to two of its vertices, and it is a face of $K_A$ since all six vertices are in $A$. For a Petrie cycle the two omitted antipodes are adjacent to three cycle vertices each, and by CP (six vertices of a commuting cube in $A$ force all eight) they belong to $A$; so in $G_A$ the cycle has an outside vertex adjacent to two of its vertices. Hence the stated criterion identifies the hexagon faces. Petrie cycles are boundaries of the three squares at an omitted antipode, so including or excluding them as $2$-cells does not change the span of boundaries. Cellular homology in degrees $\le1$ depends only on the $2$-skeleton, hence $\beta_0,\beta_1$.
\end{proof}
\begin{proposition}[negative control: a Petrie cycle and a braid hexagon]\label{prop:petrie}
In $P_6$, let $A_1$ be the six vertices of a commuting cube $\{1,s_1\}\times\{1,s_3\}\times\{1,s_5\}$ other than $x$ and $xs_1s_3s_5$, and let $A_2$ be a braid hexagon. Then $G_{A_1}\cong G_{A_2}\cong C_6$ but $K_{A_1}$ has $f=(6,6)$, $\beta=(1,1)$ while $K_{A_2}$ has $f=(6,6,1)$, $\beta=(1,0,0)$. Without CP, the uncoloured graph does not determine $\beta_1$.
\status{\FiniteExact; re-derived by the manuscript verifier (\texttt{ch5\_petrie\_vs\_braid})}{GIG-I v002, Z1; GIG-I v003, T20e}
\end{proposition}
\Cref{cor:CP} shows that no prefix-defined set has a Petrie cycle as trace of a cube without its antipodes, which is exactly why graph-only recovery works in the arithmetic class.

\section{Recognising higher cells}
Beyond dimension two, cells are recognised by convexity.
\begin{theorem}[recognition of higher cells]\label{thm:face-recognition}
(a) \emph{Completeness (all $k$, all $A$).} If $F$ is a face of $P_k$ of dimension $\ge3$ with $V(F)\subseteq A$, then the induced subgraph $G_A[V(F)]$ is the face graph $\prod_i\mathrm{Cay}(S_{|B_i|})$ and $V(F)$ is geodesically convex in $G_A$.
(b) \emph{Soundness in the full permutohedron ($k\le7$ all face types of dimension $\ge3$; $k=8$ the three types of $3$-faces).} In the graph of $P_k$ itself the convex induced copies of a $d$-face graph, $d\ge3$, are exactly the $d$-faces. Consequently graph-only reconstruction is exact for every $A$ that is geodesically convex in $P_k$ ($k\le7$).
(c) \emph{Verified domains.} For the exhaustive domains $k=4$ (labels $\le40$), $k=5$ (labels $\le16$), a stated $k=6$ sample, the two anchors, and targeted $k=7$ regimes, the face poset derived from the uncoloured graph equals the true face poset; the graph determines $\operatorname{rank}\partial_3$.
(d) \emph{Open.} For a non-convex realizable cube-product set a subset convex in the sparser graph $G_A$ and isomorphic to a face graph might not be a face; soundness there is not proved.
\status{(a) \Analytic; (b) \CompAss\ (E12 enumeration of induced copies through a root; inherited, not re-run); (c) \FiniteExact\ (inherited); (d) \textbf{open}}{GIG-I v003, T20a--d}
\end{theorem}
\begin{proof}[Proof of (a)]
Two vertices of $F$ differ only by exchanges inside blocks, which are adjacent positions, so $G_A[V(F)]$ is the product of Cayley graphs. Graph distance in $P_k$ is the Kendall-tau distance (number of label pairs ordered differently), and a path is geodesic iff each step swaps a pair that is inverted relative to the target. For $u,v\in V(F)$ only pairs inside a block are inverted, and such pairs are adjacent at the moment of swapping, so every geodesic stays in $V(F)\subseteq A$. Hence $d_A(u,v)=d_P(u,v)$ and every $G_A$-geodesic is a $P_k$-geodesic, hence lies in $V(F)$.
\end{proof}
\begin{theorem}[third-boundary obstruction]\label{thm:third-boundary}
Let $X,Y$ be realizable distinct-label complexes ($k\le8$) with isomorphic uncoloured swap graphs. Then $f_0,f_1,f_2,\rho_1,\rho_2$ coincide and $\beta_2(X)-\beta_2(Y)=\rho_3(Y)-\rho_3(X)$. Thus a pair of graph-isomorphic realizable systems with different $\beta_2$ must differ in $\operatorname{rank}\partial_3$ ($\rho_3=J_3$), and equal $\rho_3$ forces equal $\beta_2$.
\status{\Analytic given \cref{thm:skeleton}}{GIG-I v003, T19}
\end{theorem}
\begin{proof}
By \cref{thm:skeleton} $f_2,\rho_2$ are graph invariants, and $\beta_2=f_2-\rho_2-\rho_3$.
\end{proof}
Whether graph-isomorphic realizable pairs with different homology exist at all is \emph{open} (it reduces to \cref{thm:face-recognition}(d)); none was found in the verified domains.
\begin{figure}[ht]\centering\includegraphics[width=\textwidth]{figures/fig2_swap_geometry.pdf}
\caption{Legal-swap geometry. (a) The permutohedron $P_4$ (truncated octahedron, orthographic projection of the exact vertex coordinates): the $16$ admissible histories of content $\{2,3,14,15\}$ at source $210$, their $20$ legal swaps and the $4$ square faces of $K_A$; the loop passing through the hexagonal region is not filled, so $\beta=(1,1,0,0)$. (b) The two kinds of $2$-cells and a commuting cube $Q_3$ with a Petrie $6$-cycle (red), which cannot be the trace of a prefix-defined $A$ (\cref{cor:CP}).}\label{fig:swaps}
\end{figure}
